% ════════════════════════════════════════════════════════════════════════════
\part*{TERCERA PARTE. Medición del Software}
\addcontentsline{toc}{part}{TERCERA PARTE. Medición del Software}
% ════════════════════════════════════════════════════════════════════════════

\section{Introducción a la medición}
La medición del software es la actividad de la ingeniería de software que asigna
valores numéricos a atributos del producto, del proceso y de los recursos, con el fin
de comprender, controlar y mejorar la calidad \parencite{pressman2014,iso25010}. En
\textbf{SmartComedor} la medición permite cuantificar el tamaño y la estructura del
código (atributos internos) y la calidad percibida durante su uso (atributos externos),
y así sustentar decisiones de mantenimiento y mejora continua.

\subsection{Objetivos}
\begin{itemize}[leftmargin=1.4em]
  \item Cuantificar el tamaño del código fuente y su distribución por módulos.
  \item Caracterizar la estructura orientada a clases (modelos y controladores).
  \item Evaluar atributos externos de calidad según \textcite{iso25010}
        (ISO/IEC 25010): usabilidad, seguridad, portabilidad, mantenibilidad,
        rendimiento, fiabilidad, compatibilidad y adecuación funcional.
  \item Identificar valores fuera de umbral que orienten acciones de mejora.
\end{itemize}

\subsection{Alcance}
La medición abarca el código de aplicación del servidor
(\texttt{server/src}) y del cliente (\texttt{client/src}), junto con la suite de
pruebas (\texttt{server/src/tests}). Se excluyen dependencias de terceros
(\texttt{node\_modules}), artefactos de compilación y de cobertura.

\subsection{Tipos y herramientas de métricas}
Se emplearon herramientas estándar de la industria: \texttt{cloc} para el conteo de
líneas, \texttt{Jest}/\texttt{Istanbul} para cobertura, y la inspección estática del
código (modelos Sequelize, controladores y rutas). La Tabla~\ref{tab:herramientas-med}
resume las herramientas por tipo de métrica.

\begin{table}[H]\centering
\caption{Herramientas por tipo de métrica}\label{tab:herramientas-med}
\begin{tabularx}{\textwidth}{L{0.32\textwidth} X}
\toprule \textbf{Tipo de métrica} & \textbf{Herramienta / fuente} \\ \midrule
Tamaño del código (LOC) & \texttt{cloc} \\
Cobertura de pruebas & \texttt{Jest} + \texttt{Istanbul} \\
Complejidad / caminos & Análisis de grafo de flujo \parencite{mccabe1976} \\
Estructura de clases & Inspección de modelos/controladores \\
Seguridad & Lista OWASP Top 10 \parencite{owasp2021} \\
Rendimiento & Escenarios de carga (Apache JMeter) \parencite{jmeter} \\
Usabilidad & Heurísticas de Nielsen \parencite{nielsen1994} \\
\bottomrule
\end{tabularx}
\end{table}

% ────────────────────────────────────────────────────────────────────────────
\section{Medición del software — atributos internos}

\subsection{Métricas de tamaño del código fuente}
La Tabla~\ref{tab:loc-global} presenta el tamaño del código de aplicación medido con
\texttt{cloc}. El servidor y el cliente suman \textbf{14.482} líneas de código efectivo
(LOC) en TypeScript, con una suite de pruebas de \textbf{3.194} LOC.

\begin{table}[H]\centering
\caption{Tamaño del código de aplicación (cloc)}\label{tab:loc-global}
\begin{tabularx}{\textwidth}{X c c c c}
\toprule
\textbf{Ámbito} & \textbf{Archivos} & \textbf{Código} & \textbf{Comentario} & \textbf{Blanco} \\
\midrule
Servidor (\texttt{server/src}) & 55 & 7.336 & 770 & 1.398 \\
\quad de los cuales, pruebas & 10 & 3.194 & --- & --- \\
Cliente (\texttt{client/src}) & 43 & 7.149 & 6 & 483 \\
\textbf{Total aplicación} & \textbf{98} & \textbf{14.485} & \textbf{776} & \textbf{1.881} \\
\bottomrule
\end{tabularx}
\end{table}

La Tabla~\ref{tab:loc-modulos} desglosa el código del servidor por capa, evidenciando
que los controladores (lógica de negocio) concentran el mayor volumen, seguidos de los
modelos y los servicios.

\begin{table}[H]\centering
\caption{Tamaño por capa del servidor (cloc)}\label{tab:loc-modulos}
\begin{tabularx}{\textwidth}{X c c}
\toprule
\textbf{Capa (\texttt{server/src})} & \textbf{Archivos} & \textbf{LOC} \\
\midrule
Controladores (\texttt{controllers}) & 9 & 1.805 \\
Modelos (\texttt{models}) & 12 & 1.096 \\
Servicios (\texttt{services}) & 5 & 342 \\
Rutas (\texttt{routes}) & 11 & 267 \\
Middleware (\texttt{middleware}) & 3 & 152 \\
Pruebas (\texttt{tests}) & 10 & 3.194 \\
\bottomrule
\end{tabularx}
\end{table}

Un indicador derivado relevante es la \textbf{razón de código de prueba a código de
producción} del servidor: $3.194 / (7.336-3.194) \approx 0{,}77$, es decir, casi
0,8 líneas de prueba por cada línea de producción, lo que refleja una inversión alta en
verificación. La densidad de comentarios del servidor es
$770/7.336 \approx 10{,}5\%$.

\subsection{Métricas de tamaño orientadas a clases}
El modelo de dominio se compone de \textbf{11 entidades} (clases Sequelize) y la lógica
de negocio se organiza en \textbf{9 controladores} que exponen \textbf{56 endpoints}
funcionales. La Tabla~\ref{tab:clases-tam} resume las métricas orientadas a clases.

\begin{table}[H]\centering
\caption{Métricas de tamaño orientadas a clases}\label{tab:clases-tam}
\begin{tabularx}{\textwidth}{X c}
\toprule \textbf{Métrica} & \textbf{Valor} \\ \midrule
Número de clases de dominio (modelos) & 11 \\
Número de controladores & 9 \\
Número de servicios & 5 \\
Endpoints REST funcionales & 56 \\
Promedio de LOC por modelo & $1.096/11 \approx 99{,}6$ \\
Promedio de LOC por controlador & $1.805/9 \approx 200{,}6$ \\
Promedio de endpoints por controlador & $56/9 \approx 6{,}2$ \\
\bottomrule
\end{tabularx}
\end{table}

Las 11 entidades del dominio son: \texttt{User}, \texttt{Student}, \texttt{Payment},
\texttt{MealAttendance}, \texttt{Rating}, \texttt{Complaint}, \texttt{News},
\texttt{Cycle}, \texttt{AuditLog}, \texttt{SupervisorAssignment} y
\texttt{SupervisorLog}. El controlador de mayor tamaño es \texttt{authController.ts}, lo
que es coherente con la cantidad de responsabilidades de seguridad (login, refresh,
recuperación de contraseña y 2FA) y con sus 59 casos de prueba.

% ────────────────────────────────────────────────────────────────────────────
\section{Medición de la calidad — atributos externos}
Se evalúan los ocho atributos externos pertinentes del modelo de calidad de
\textcite{iso25010}. Para cada uno se define una métrica, su fórmula y el valor
observado en SmartComedor.

\subsection{Métricas de usabilidad}
Evaluadas con las heurísticas de \textcite{nielsen1994} sobre la interfaz Material
Design~3 \parencite{material3}. La Tabla~\ref{tab:m-usab} resume la inspección.

\begin{table}[H]\centering
\caption{Métricas de usabilidad (heurísticas de Nielsen)}\label{tab:m-usab}
\begin{tabularx}{\textwidth}{L{0.46\textwidth} C{0.16\textwidth} X}
\toprule \textbf{Heurística} & \textbf{Cumplimiento} & \textbf{Evidencia} \\ \midrule
Visibilidad del estado del sistema & Alto & Spinners, snackbars y estados de carga \\
Correspondencia con el mundo real & Alto & Lenguaje del dominio (almuerzo, ciclo) \\
Control y libertad del usuario & Medio & Acciones reversibles en formularios \\
Consistencia y estándares & Alto & Componentes MUI reutilizados \\
Prevención de errores & Alto & Validación en formularios y listas cerradas \\
Reconocer antes que recordar & Alto & Búsqueda por nombre/UID en supervisión \\
Flexibilidad y eficiencia & Medio & Atajos por rol; exportación configurable \\
Diseño estético y minimalista & Alto & Jerarquía visual M3 \\
Ayuda ante errores & Medio & Mensajes guiados de validación \\
Ayuda y documentación & Medio & Textos de ayuda contextual \\
\bottomrule
\end{tabularx}
\end{table}
\noindent\textbf{Índice de usabilidad heurística:} 7 heurísticas en nivel Alto y 3 en
nivel Medio sobre 10 evaluadas, lo que arroja un cumplimiento ponderado aproximado del
\textbf{85\%}.

\subsection{Métricas de seguridad}
Medidas contra OWASP Top~10 \parencite{owasp2021}. La Tabla~\ref{tab:m-seg} indica el
control implementado por categoría de riesgo.

\begin{table}[H]\centering
\caption{Cobertura de controles de seguridad (OWASP Top 10:2021)}\label{tab:m-seg}
\begin{tabularx}{\textwidth}{L{0.40\textwidth} C{0.14\textwidth} X}
\toprule \textbf{Riesgo} & \textbf{Estado} & \textbf{Control} \\ \midrule
A01 Control de acceso roto & Mitigado & Middleware JWT + autorización por rol \\
A02 Fallos criptográficos & Mitigado & Hash bcrypt; secretos por entorno \\
A03 Inyección & Mitigado & ORM Sequelize (consultas parametrizadas) \\
A04 Diseño inseguro & Mitigado & Validación de reglas de negocio \\
A05 Mala configuración & Parcial & Variables de entorno; CORS controlado \\
A07 Fallas de identificación y autenticación & Mitigado & 2FA TOTP; expiración de tokens \\
A09 Fallos de registro y monitoreo & Mitigado & \texttt{audit\_logs} y \texttt{supervisor\_logs} \\
\bottomrule
\end{tabularx}
\end{table}
\noindent\textbf{Métrica de cobertura de seguridad de la lógica de acceso:}
la cobertura de pruebas del \texttt{middleware/auth.ts} es del \textbf{97,5\%} de
sentencias y \textbf{91,66\%} de ramas, valores apropiados para un componente crítico.

\subsection{Métricas de portabilidad}
La portabilidad se cuantifica como la capacidad de ejecutar el sistema en distintos
entornos sin cambios de código. SmartComedor se distribuye con Docker, lo que produce
un índice de portabilidad alto (Tabla~\ref{tab:m-port}).

\begin{table}[H]\centering
\caption{Métricas de portabilidad}\label{tab:m-port}
\begin{tabularx}{\textwidth}{L{0.46\textwidth} X}
\toprule \textbf{Métrica} & \textbf{Valor} \\ \midrule
Adaptabilidad (entornos soportados) & Docker (\texttt{dev}/\texttt{prod}) y ejecución local \\
Facilidad de instalación & \texttt{docker compose --profile dev up} (un comando) \\
Independencia de plataforma & Node.js + contenedores (Linux/macOS/Windows) \\
Navegadores objetivo & Chrome, Firefox, Edge, Safari (SPA responsive) \\
\bottomrule
\end{tabularx}
\end{table}

\subsection{Métricas de mantenibilidad}
La mantenibilidad se estima con la modularidad, la analizabilidad (cobertura) y la
capacidad de prueba. La Tabla~\ref{tab:m-mant} resume los indicadores.

\begin{table}[H]\centering
\caption{Métricas de mantenibilidad}\label{tab:m-mant}
\begin{tabularx}{\textwidth}{L{0.46\textwidth} X}
\toprule \textbf{Métrica} & \textbf{Valor} \\ \midrule
Modularidad & Separación en capas (rutas/controladores/servicios/modelos) \\
Densidad de comentarios (servidor) & $\approx 10{,}5\%$ \\
Razón prueba/producción & $\approx 0{,}77$ \\
Analizabilidad (cobertura global) & 80,44\% sentencias \\
Capacidad de prueba & 196 casos automatizados, deterministas \\
\bottomrule
\end{tabularx}
\end{table}

\subsection{Métricas de rendimiento}
El objetivo de rendimiento es que las operaciones críticas (autenticación y consultas)
respondan en menos de 3~segundos \parencite{jmeter}. La Tabla~\ref{tab:m-rend} resume
los indicadores de eficiencia de desempeño.

\begin{table}[H]\centering
\caption{Métricas de rendimiento}\label{tab:m-rend}
\begin{tabularx}{\textwidth}{L{0.46\textwidth} X}
\toprule \textbf{Métrica} & \textbf{Objetivo / valor} \\ \midrule
Tiempo de respuesta (operaciones críticas) & $< 3$~s \\
Comportamiento temporal & Consultas indexadas por clave/UID \\
Utilización de recursos & Contenedores con \texttt{healthcheck} \\
Tiempo de ejecución de la suite & $\approx 4{,}7$~s (196 casos) \\
\bottomrule
\end{tabularx}
\end{table}

\subsection{Métricas de fiabilidad}
La fiabilidad se evidencia en la tasa de éxito de las pruebas y en los mecanismos de
recuperación. La Tabla~\ref{tab:m-fiab} resume los indicadores.

\begin{table}[H]\centering
\caption{Métricas de fiabilidad}\label{tab:m-fiab}
\begin{tabularx}{\textwidth}{L{0.46\textwidth} X}
\toprule \textbf{Métrica} & \textbf{Valor} \\ \midrule
Madurez (tasa de éxito de pruebas) & 100\% (196/196) \\
Densidad de defectos abiertos & 0 en la suite automatizada \\
Tolerancia a fallos & Manejo de errores con códigos HTTP controlados \\
Recuperabilidad & Respaldo diario con \texttt{pg\_dump} (retención 30 días) \\
\bottomrule
\end{tabularx}
\end{table}

\subsection{Métricas de compatibilidad}
La compatibilidad mide la coexistencia e interoperabilidad. SmartComedor expone una API
REST con JSON, consumible por el cliente web y por terceros autorizados
(Tabla~\ref{tab:m-comp}).

\begin{table}[H]\centering
\caption{Métricas de compatibilidad}\label{tab:m-comp}
\begin{tabularx}{\textwidth}{L{0.46\textwidth} X}
\toprule \textbf{Métrica} & \textbf{Valor} \\ \midrule
Interoperabilidad & API REST/JSON con autenticación JWT \\
Coexistencia & Servicios desacoplados en contenedores \\
Formatos de exportación & Excel, CSV y PDF \\
\bottomrule
\end{tabularx}
\end{table}

\subsection{Métricas de adecuación funcional}
Mide en qué grado el sistema cubre las funciones requeridas. Con 45 requisitos
funcionales y 56 endpoints implementados, la cobertura funcional es completa
(Tabla~\ref{tab:m-adec}).

\begin{table}[H]\centering
\caption{Métricas de adecuación funcional}\label{tab:m-adec}
\begin{tabularx}{\textwidth}{L{0.46\textwidth} X}
\toprule \textbf{Métrica} & \textbf{Valor} \\ \midrule
Completitud funcional & 45/45 requisitos funcionales implementados \\
Corrección funcional & Verificada por 196 casos de prueba \\
Pertinencia funcional & 56 endpoints alineados a los casos de uso \\
\bottomrule
\end{tabularx}
\end{table}

\section{Conclusiones de la medición}
La medición confirma un producto de tamaño medio ($\approx 14.5$~KLOC de aplicación)
con una fuerte inversión en pruebas (razón prueba/producción $\approx 0{,}77$ y
cobertura global del 80,44\%). Los atributos externos más críticos —seguridad y
fiabilidad— presentan los mejores indicadores (cobertura del componente de acceso
$\geq 97\%$ y 100\% de éxito en pruebas). Las oportunidades de mejora se concentran en
elevar la cobertura de \texttt{mealsController.ts} y \texttt{auditService.ts}, y en
formalizar mediciones continuas de rendimiento en producción.
